\begin{helper}[GnpCodegreeBound]
\label{helper:GnpCodegreeBound}
Let \(m\) and \(T\) be natural numbers. Let \(0 \le p \le 1\).
Let \(r\) and \(s\) be distinct vertices in \(\operatorname{Fin}(m)\).
Under the \(G(m,p)\) product law, the probability that \(r\) and \(s\) have at least \(T\) common neighbors is bounded by
\[
\binom{m}{T} (p^2)^T.
\]
Explicitly,
\[
\sum_{G : T \le |\{t : G.Adj(r,t) \land G.Adj(s,t)\}|} \noderef{GnpGraphWeight}(p, G) \le \binom{m}{T} (p^2)^T.
\]
\end{helper}

\begin{proof}
Expand the graph probability using \(\noderef{GnpGraphWeight}\).  Its weight
is a product over unordered edge coordinates: a present edge contributes
\(p\), and an absent edge contributes \(1-p\).  Since \(0\le p\le1\), every
factor and hence every graph weight is nonnegative.

Let
\[
C_G=\{t\in\operatorname{Fin}(m):G.Adj(r,t)\land G.Adj(s,t)\}.
\]
If \(|C_G|\ge T\), then there is a \(T\)-element subset
\(S\subseteq C_G\).  Thus the event in the statement is contained in the
union, over all \(T\)-subsets \(S\subseteq\operatorname{Fin}(m)\), of the
cylinder event
\[
E_S=\{G:\forall t\in S,\ G.Adj(r,t)\land G.Adj(s,t)\}.
\tag{1}
\]
Because the graph weights are nonnegative, the finite union bound gives
\[
\sum_{G:|C_G|\ge T}\noderef{GnpGraphWeight}(p,G)
\le
\sum_{S:\ |S|=T}\ \sum_{G\in E_S}\noderef{GnpGraphWeight}(p,G).
\tag{2}
\]

Fix such an \(S\).  If \(r\in S\), then \(G.Adj(r,r)\) would be required in
\(E_S\), impossible by looplessness of a simple graph; similarly \(s\in S\)
makes \(E_S\) empty.  In either case the inner sum is \(0\), hence at most
\((p^2)^T\).

Assume now that \(S\) contains neither \(r\) nor \(s\).  For every \(t\in S\),
the cylinder \(E_S\) forces the two distinct edge coordinates
\(\{r,t\}\) and \(\{s,t\}\) to be present.  As \(t\) varies in \(S\), these
\(2T\) edge coordinates are all distinct because \(r\ne s\) and
\(t\notin\{r,s\}\).  Summing \(\noderef{GnpGraphWeight}\) over \(E_S\), the
forced coordinates contribute \(p^{2T}=(p^2)^T\).  Every unforced edge
coordinate is summed over its two possibilities and contributes
\[
p+(1-p)=1.
\]
Therefore
\[
\sum_{G\in E_S}\noderef{GnpGraphWeight}(p,G)=(p^2)^T.
\tag{3}
\]
Combining the empty-case estimate and (3), every \(T\)-subset contributes at
most \((p^2)^T\).  There are \(\binom mT\) such subsets.  Applying this to
(2) proves
\[
\sum_{G:|C_G|\ge T}\noderef{GnpGraphWeight}(p,G)
\le
\binom mT(p^2)^T,
\]
which is the claimed bound.
\end{proof}
